\chapter{Various}

\section{Random}
\kactlimport{RNG.h}

\section{Ternary Search}
\kactlimport{GoldenSectionSearch.h}
\kactlimport{TernarySearch.h}

\section{Optimization tricks}
	\subsection{Bit hacks}
		\begin{itemize}
			\item \verb@x & -x@ is the least bit in \texttt{x}.
			\item \verb@for (int x = m; x; ) { --x &= m; ... }@ loops over all subset masks of \texttt{m} (except \texttt{m} itself).
			\item \verb@c = x&-x, r = x+c; (((r^x) >> 2)/c) | r@ is the next number after \texttt{x} with the same number of bits set.
			\item \verb@rep(b,0,K) rep(i,0,(1 << K))@ \\ \verb@  if (i & 1 << b) D[i] += D[i^(1 << b)];@ computes all sums of subsets.
		\end{itemize}
	\subsection{Pragmas}
		\begin{itemize}
			\item \lstinline{#pragma GCC optimize("Ofast")} will make GCC auto-vectorize loops and aggressively optimize floating points.
			\item \lstinline{#pragma GCC target("avx2")} can improve performance of vectorized code, but causes crashes on old machines.
			\item \lstinline{#pragma GCC target("bmi,bmi2,lzcnt,popcnt")} are good for bitsets.
		\end{itemize}
